\section{Phần mở đầu}
%\addcontentsline{toc}{part}{Phần mở đầu}
\fontsize{12pt}{0pt}
\selectfont
\subsection{Lý do chọn đề tài}
Trong bối cảnh kinh tế hiện đại, các cửa hàng bán lẻ, đặc biệt là cửa hàng điện tử, đối mặt với sự cạnh tranh gay gắt và biến động không ngừng của thị trường. Để tồn tại và phát triển, các cửa hàng cần không chỉ hiểu rõ nhu cầu của khách hàng mà còn tối ưu hóa mọi khía cạnh trong hoạt động kinh doanh, từ quản lý hàng tồn kho, lên kế hoạch bán hàng, đến đánh giá hiệu quả của các chiến dịch tiếp thị.

Trên cơ sở đó, đề tài "Thống kê dữ liệu về doanh số bán lẻ của một cửa hàng điện tử" được nhóm tác giả lựa chọn nhằm nghiên cứu sâu hơn về cách mà dữ liệu doanh số có thể cung cấp thông tin quan trọng cho việc ra quyết định kinh doanh. Thông qua phân tích thống kê dữ liệu bán hàng, cửa hàng có thể nhận diện các xu hướng tiêu dùng và đánh giá hiệu quả các hoạt động khuyến mãi. Điều này giúp cung cấp cái nhìn toàn diện về sức khỏe tài chính của cửa hàng.
\subsection{Mục tiêu nghiên cứu}
Mục tiêu của bài tập lớn này là tìm hiểu sâu hơn về các nguyên tắc cơ bản và ứng dụng của lý thuyết thống kê dựa trên dữ liệu về doanh số bán lẻ của một cửa hàng điện tử. Bên cạnh đó, việc nghiên cứu đề tài còn giúp chúng ta cải thiện kỹ năng đọc dữ liệu, làm sạch dữ liệu và phân tích dữ liệu vào thực tiễn cũng như cải thiện kĩ năng tư duy logic.
\subsection{Nhiệm vụ nghiên cứu}
Đầu tiên, ta sẽ nhập dữ liệu vào môi trường R từ nguồn cung cấp, có thể từ file CSV và file Excel . Mục tiêu là ta làm quen với cấu trúc của dữ liệu và đảm bảo dữ liệu đã được tải đúng và đầy đủ.

Tiếp theo, những dữ liệu khuyết sẽ được xử lý (NA) hoặc các giá trị không hợp lệ, đảm bảo tính chính xác và sẵn sàng của dữ liệu cho phân tích.

Các biến dữ liệu sẽ được chuyển đổi (như log-transform, chuẩn hóa hoặc chuẩn hóa lại loại dữ liệu) để đảm bảo phù hợp với các phân tích thống kê tiếp theo.

Hơn nữa, Thống kê mô tả sẽ được thực hiện để hiểu rõ các đặc điểm cơ bản của dữ liệu. Chúng tôi sẽ sử dụng cả các thống kê mẫu (như trung bình, độ lệch chuẩn) và biểu đồ để trình bày các thông tin quan trọng từ dữ liệu.

Toàn bộ các bước thực hiện trong nghiên cứu sẽ được nhóm tác giả triển khai bằng ngôn ngữ R. Các đoạn mã sẽ được chú thích cẩn thận, giúp minh bạch hóa quá trình phân tích, cũng như đảm bảo khả năng tái lập kết quả cho những người đọc có nhu cầu tìm hiểu sâu hơn.
\subsection{Phạm vi nghiên cứu}
Phạm vi về mặt thời gian: Đề tài tập trung nghiên cứu về hoạt động bán hàng trong năm 2019. Khoảng thời gian một năm được lựa chọn nhằm cung cấp cái nhìn tổng quát và chi tiết về xu hướng doanh số, bao gồm các giai đoạn cao điểm như dịp lễ Tết và các chương trình khuyến mãi lớn trong năm.
